\section{Design Consequences and Limits}\label{sec:consequences}

\subsection{The three contracts compared}\label{cr:compare}


\Cref{tab:contracts} puts the floor-cap contract beside the two extensions, with numbers for the linear price $p(q)=1+q$, a total lock of one token, loan caps of one and no fees (script block N17). The buffered contract adds $40\%$ to the credit of every lock for a reserve of $0.08$, and exposes that reserve in full. The curve-valued contract needs no reserve and has no such exposure, but what it adds depends on the size of a lock: half the lock size in relative terms on this curve, so $5\%$ for a lock of $0.1$ and $50\%$ for a single lock of the whole supply, and nothing at all for a lock that fits inside a flat floor segment. Ten equal locks together obtain $5\%$, one pooled lock $50\%$; the difference is the price of independent exit. Whether either extension is worth its cost depends on the distribution of lock sizes and on what the reserve could otherwise earn, which we do not measure.

\begin{table}[t]
\centering\small
\caption{The floor-cap contract and the two extensions of this section. Numbers are for $p(q)=1+q$, total lock $L=1$ (floor value $1$), caps $\gamma=\gbar=1$ and no fees; ``gain'' is the coalition gain of \cref{cr:extraction}.}\label{tab:contracts}
\begin{tabularx}{\textwidth}{@{}>{\raggedright\arraybackslash}p{0.2\textwidth}>{\raggedright\arraybackslash}X>{\raggedright\arraybackslash}X>{\raggedright\arraybackslash}X@{}}
\toprule
 & Floor cap (\cref{sec:model}) & Committed reserve, proportional release & Curve-valued lock\\
\midrule
Credit of a lock $\ell$ & $\Pf\ell=\ell$ & $(D_{\max}/L)\,\ell=1.4\,\ell$ at $W=0.08$ & $R(\ell)=\ell+\ell^2/2$\\
Aggregate credit & $1$ & $1.4$ & $1.05$ (ten locks) to $1.5$ (one lock)\\
Committed capital & none & $W=0.08$ & none\\
Largest gain & $0$ & $W_{\min}=0.08$, the whole reserve & $0$\\
Release on partial repayment & proportional & proportional & less than proportional at the limit: $0.419$ of a unit lock for half the debt\\
Splitting a position into independent loans & leaves credit unchanged & leaves credit unchanged & reduces capacity; unrestricted proportional splitting forces the floor bound (\cref{cr:splitting})\\
Guards on curve changes & (G1), (G2) & reserve test with the new curve & (G1), (G2) and the loan condition under the new curve\\
Check at a borrow & one inequality per loan & sorted envelope of the whole book & one inequality per loan, with $R$\\
Evidence in this paper & Lean proof, reference model, contract, tests & proofs, exact finite checks & proofs, exact finite checks\\
\bottomrule
\end{tabularx}
\end{table}

\subsection{Who funds the reserve}

\Cref{cr:extraction} changes how a committed reserve should be read. It protects those who redeem. Under proportional release it is also what a coalition holding all the tokens \emph{can} take: with no fees and a reserve funded by others, at most $W_{\min}$, and all of $W$ at the maximum admitted credit. With fees, the gain is reduced by the amounts in part~(d) of the proposition and can be negative. We have proved that the strategy is available and what it pays at most; we have not valued the option that it gives borrowers, and we do not predict that it is exercised. A sponsor who funds a reserve to offer more credit should nevertheless count it as capital exposed to that strategy and not as idle liquidity. If the coalition funds the reserve itself, the contribution has to be subtracted from its proceeds, and under the same assumptions its net gain is at most $W_{\min}-W\le0$; the reserve still does its work of covering redemptions, and it can still move value between borrowers who contribute and repay unequally. The curve-valued contract avoids the exposure by never lending more than the reserve value of what stays locked, and pays for that with the restrictions of \cref{cr:pooling}.

\subsection{What safety does not give}

The results above are statements about a balance sheet in units of the reserve asset. They do not show that anyone wants the credit. The extended version \citep{extended} analyzes the floor-cap contract further and finds that safety and viability separate there too. For a price-taking borrower who values cash flows, holds tokens that pay nothing and has no use for them in the meantime, deferring the repayment of a loan without carrying cost is weakly optimal. And in a stationary equilibrium at the floor, with given arrivals and no retained surplus, entrants as a group pay for the reserve that a rising floor needs and for the time their capital spends in the system: the token-weighted mean value of entering is negative whenever the floor grows more slowly than holders discount and the loan cap is below the net floor. In event data from one deployed system of this kind, $\ELeverAmtPct\%$ of new debt by amount finances purchases of more of the same token. A theory of that leveraged demand is the missing piece, and neither version of this paper supplies it.

Four further limits apply. The floor never falls in units of the reserve asset, which says nothing about the value of that asset. The capacity results fix the curve; a design that changes the curve needs the guards of \cref{cr:guards} or the stronger reshape guard of \cref{cr:nonlinear}. The alternative contracts are specified and proved on paper, and are neither implemented nor machine-checked. And we have not measured how much credit the curve-valued contract would add at the lock sizes of a real system; for a lock that fits inside a flat floor segment it adds nothing.

\paragraph{Disclosure.} The author is a co-founder of House Markets, which develops token issuance and credit mechanisms of the kind analyzed here, and is a co-author of \vone, on which the floor-cap contract builds. The author holds no tokens of Baseline or Nirvana.

\section{Conclusion}

Lending the redemption reserve against locked tokens is safe today whenever debt is covered by the reserve of the locked part of the curve. It stays safe only if residual debt is covered by the committed reserve plus the reserve of the residual lock in every state that borrowers can reach by repaying, and this is a condition on the withdrawal rights that the loans carry. With proportional release and no committed reserve, the condition is the floor cap, loan by loan. Lending more needs a reserve whose exact size is the largest gap between a sorted debt envelope and the reserve curve, and with no fees that same number is the most that a coalition holding all the tokens can take out of a reserve funded by others. A contract that releases collateral along the curve lends more without a reserve and without that exposure, and gives up proportional release, free splitting and the simple reshape guards. More credit therefore has a price in capital or in withdrawal rights, and the choice between the two is a design decision that the accounting makes explicit.
